\chapter[Band 19: Reelle Zahlen]{Band 19 auf einen Blick}
\noindent\textbf{Reelle Zahlen}\par\medskip
\noindent\emph{Lesart.}
\(A,B,C\) sind reelle Zahlen, im Modell also Dedekindsche Schnitte;
\(p,q,r\) sind rationale Zahlen. In dieser Übersicht stehen
\(0,1,+,-,\cdot,\leq,<\) bei reellen Argumenten für die im Band konstruierten
Schnittoperationen und die Schnittordnung. Rationale Operationen innerhalb
einer Schnittbeschreibung bleiben rationale Operationen.

\begin{BandUebersicht}
\textbf{Dedekindsche Schnitte und reeller Träger}
\(\RealCut A\) wird durch die Schnittaxiome festgelegt:
\UebFormel{\begin{aligned}
&\varnothing\neq A\subsetneq\mathbb Q,\\
&p\in A,\ q\in\mathbb Q,\ q<p\vdash q\in A,\\
&p\in A\vdash\exists r\in A\;(p<r),\\
&\mathbb R\coloneqq
\{A\in\powerset(\mathbb Q)\mid\RealCut A\}.
\end{aligned}}
\UebQuellen{\FormulaRefAuto{DedekindCutPredicate}[def];
\FormulaRefAuto{DedekindCutNonemptyAxiom}[ax];
\FormulaRefAuto{DedekindCutDownwardAxiom}[ax];
\FormulaRefAuto{DedekindCutNoMaximumAxiom}[ax];
\FormulaRefAuto{RealCutCarrier}[def]}
&
\textit{Schnittkriterium und Elementcharakterisierung}
Genau die links genannten drei Bedingungen an \(A\) sind
äquivalent zur Schnitteigenschaft.
\UebFormel{A\in\mathbb R\eqvdash\RealCut A.}
Damit werden die axiomatisch beschriebenen Schnitte als Elemente
einer Menge zusammengefasst.
\UebQuellen{\FormulaRefAuto{DedekindCutCharacterization}[thm];
\FormulaRefAuto{RealCutMembership}[thm]}
\\ \addlinespace[.8ex]

\textbf{Hauptschnitt und rationale Einbettung}
\UebFormel{\begin{aligned}
\RealHat q&\coloneqq\{p\in\mathbb Q\mid p<q\},\\
\RealEmb&\colon\mathbb Q\to\mathbb R,\\
\RealEmb(q)&\coloneqq\RealHat q.
\end{aligned}}
\UebQuellen{\FormulaRefAuto{RealPrincipalCut}[def];
\FormulaRefAuto{RationalToRealEmbedding}[def]}
&
\textit{Injektive und ordnungstreue Einbettung}
\UebFormel{\begin{aligned}
&\RealCut{\RealHat q},\qquad
\RealEmb\colon\mathbb Q\inj\mathbb R,\\
&p\leq q\leftrightarrow\RealEmb(p)\leq\RealEmb(q),\\
&p<q\leftrightarrow\RealEmb(p)<\RealEmb(q).
\end{aligned}}
\UebQuellen{\FormulaRefAuto{RealPrincipalCutIsCut}[thm];
\FormulaRefAuto{RationalToRealEmbeddingInjective}[thm];
\FormulaRefAuto{RationalToRealEmbeddingOrder}[thm]}
\\ \addlinespace[.8ex]

\textbf{Schnittordnung}
\UebFormel{\begin{aligned}
A\RealLe B&\coloneqq A\subseteq B,\\
A\RealLt B&\coloneqq A\subsetneq B.
\end{aligned}}
Für die Vollständigkeit betrachtet man
\(\varnothing\neq\mathcal A\subseteq\mathbb R\) mit einer oberen
Schranke \(U\in\mathbb R\).
\UebQuellen{\FormulaRefAuto{RealCutOrder}[def];
\FormulaRefAuto{RealCutStrictOrder}[def]}
&
\textit{Totale und vollständige Ordnung}
\UebFormel{\begin{aligned}
&\TotOrd{\mathbb R,\RealLe},\\
&(\forall A\in\mathcal A\;A\leq U)\vdash\\
&\quad\bigcup\mathcal A\in\mathbb R,\qquad
\sup_{\mathbb R,\RealLe}\mathcal A=\bigcup\mathcal A.
\end{aligned}}
Eine nichtleere nach unten beschränkte reelle Menge besitzt
entsprechend ein reelles Infimum.
\UebQuellen{\FormulaRefAuto{RealCutTotalOrder}[thm];
\FormulaRefAuto{RealLeastUpperBoundProperty}[thm];
\FormulaRefAuto{RealCutSupremumUnion}[thm];
\FormulaRefAuto{RealGreatestLowerBoundProperty}[thm]}
\\ \addlinespace[.8ex]

\textbf{Kanonische Zahlkopien}
\UebFormel{\begin{aligned}
\RatCopyInReal&\coloneqq\RealEmb[\mathbb Q],\\
\IntCopyInReal&\coloneqq\RealEmb[\IntCopyInRat],\\
\NatCopyInReal&\coloneqq\RealEmb[\NatCopyInRat].
\end{aligned}}
\UebQuellen{\FormulaRefAuto{RationalCopyInReals}[def];
\FormulaRefAuto{IntegerCopyInReals}[def];
\FormulaRefAuto{NaturalCopyInReals}[def]}
&
\textit{Geschachtelter Zahlbereichsturm}
\UebFormel{\begin{aligned}
&\NatCopyInReal\subseteq\IntCopyInReal,\\
&\IntCopyInReal\subseteq\RatCopyInReal
\subseteq\mathbb R,\\
&\RealEmb\colon\mathbb Q\bij\RatCopyInReal.
\end{aligned}}
Die Identifikation erfolgt über die Einbettungen und ihre
Einschränkungen.
\UebQuellen{\FormulaRefAuto{EmbeddedNumberTowerInReals}[thm];
\FormulaRefAuto{RationalCopyRealBijection}[thm];
\FormulaRefAuto{IntegerCopyRealBijection}[thm];
\FormulaRefAuto{NaturalCopyRealBijection}[thm]}
\\ \addlinespace[.8ex]

\textbf{Addition und additive Null}
\UebFormel{\begin{aligned}
A+B\coloneqq\{q\in\mathbb Q\mid{}&
\exists a\in A\;\exists b\in B\\
&q<a+b\},\\
0&\coloneqq\RealHat{\RatZero}.
\end{aligned}}
\UebQuellen{\FormulaRefAuto{RealCutAddition}[def];
\FormulaRefAuto{RealCutZero}[def]}
&
\textit{Abschluss und additive Gesetze}
\UebFormel{\begin{aligned}
&A+B\in\mathbb R,\qquad0\in\mathbb R,\\
&(A+B)+C=A+(B+C),\\
&A+B=B+A,\qquad A+0=A.
\end{aligned}}
\UebQuellen{\FormulaRefAuto{RealCutAdditiveClosure}[thm];
\FormulaRefAuto{RealCutAdditiveLaws}[thm]}
\\ \addlinespace[.8ex]

\textbf{Negation eines Schnittes}
\UebFormel{\begin{aligned}
-A\coloneqq\{q\in\mathbb Q\mid{}&
\exists r\in\mathbb Q\\
&(r\notin A\land q<-r)\}.
\end{aligned}}
\UebQuellen{\FormulaRefAuto{RealCutNegation}[def]}
&
\textit{Additives Inverses und Ordnungsumkehr}
\UebFormel{\begin{aligned}
&-A\in\mathbb R,\qquad A+(-A)=0,\\
&-(-A)=A,\qquad -0=0,\\
&A\leq B\leftrightarrow-B\leq-A.
\end{aligned}}
\UebQuellen{\FormulaRefAuto{RealCutAdditiveClosure}[thm];
\FormulaRefAuto{RealCutAdditiveLaws}[thm];
\FormulaRefAuto{RealCutNegationOrder}[thm]}
\\ \addlinespace[.8ex]

\textbf{Positives Produkt und Vorzeichenfortsetzung}
Für \(0\leq A,B\):
\UebFormel{\begin{aligned}
A\RealPosMul B\coloneqq{}&0\cup
\{q\in\mathbb Q\mid\\
&\exists a\in A\;\exists b\in B\\
&(\RatZero<a\land\RatZero<b\land q<ab)\}.
\end{aligned}}
Das allgemeine Produkt entsteht durch die vier Vorzeichenfälle;
\(1\coloneqq\RealHat{\RatOne}\).
\UebQuellen{\FormulaRefAuto{RealCutPositiveProduct}[def];
\FormulaRefAuto{RealCutMultiplication}[def];
\FormulaRefAuto{RealCutOne}[def]}
&
\textit{Multiplikation und Ordnung}
\UebFormel{\begin{aligned}
&AB\in\mathbb R,\qquad(AB)C=A(BC),\\
&AB=BA,\qquad A1=A,\\
&A(B+C)=AB+AC,\\
&0\leq A,\ 0\leq B\vdash0\leq AB.
\end{aligned}}
\UebQuellen{\FormulaRefAuto{RealCutMultiplicativeClosure}[thm];
\FormulaRefAuto{RealCutMultiplicativeLaws}[thm];
\FormulaRefAuto{RealCutOrderArithmeticCompatibility}[thm]}
\\ \addlinespace[.8ex]

\textbf{Reziproker Schnitt}
Für \(0<A\):
\UebFormel{\begin{aligned}
&\RealPosInv A\coloneqq0\cup
\{q\in\mathbb Q\mid\\
&\quad\exists r\in\mathbb Q\;(\RatZero<r\land r\notin A\\
&\qquad{}\land q<r^{-1})\}.
\end{aligned}}
Für \(A<0\) ist \(A^{-1}\coloneqq-\RealPosInv{(-A)}\).
\UebQuellen{\FormulaRefAuto{RealCutPositiveReciprocal}[def];
\FormulaRefAuto{RealCutReciprocal}[def]}
&
\textit{Inversengesetz und Erhalt der rationalen Operationen}
\UebFormel{\begin{aligned}
&A\neq0\vdash A^{-1}\in\mathbb R,\quad AA^{-1}=1,\\
&\RealEmb(p+q)=\RealEmb(p)+\RealEmb(q),\\
&\RealEmb(pq)=\RealEmb(p)\RealEmb(q),\\
&p\neq\RatZero\vdash
\RealEmb(p^{-1})=(\RealEmb(p))^{-1}.
\end{aligned}}
\UebQuellen{\FormulaRefAuto{RealCutMultiplicativeClosure}[thm];
\FormulaRefAuto{RealCutMultiplicativeLaws}[thm];
\FormulaRefAuto{RationalToRealEmbeddingArithmetic}[thm]}
\\ \addlinespace[.8ex]

\textbf{Betrag und Abstand}
\UebFormel{\begin{aligned}
|A|&\coloneqq
\begin{cases}A,&0\leq A,\\-A,&A<0,\end{cases}\\
\RealDist A B&\coloneqq|A-B|.
\end{aligned}}
\UebQuellen{\FormulaRefAuto{RealAbsoluteValue}[def];
\FormulaRefAuto{RealDistance}[def]}
&
\textit{Betrags- und Abstandsgesetze}
\UebFormel{\begin{aligned}
&|AB|=|A||B|,\qquad|A+B|\leq|A|+|B|,\\
&0\leq\RealDist A B,\\
&\RealDist A B=0\leftrightarrow A=B,\\
&\RealDist A B=\RealDist B A,\\
&\RealDist A C\leq\RealDist A B+\RealDist B C.
\end{aligned}}
\UebQuellen{\FormulaRefAuto{RealAbsoluteValueLaws}[thm];
\FormulaRefAuto{RealDistanceLaws}[thm]}
\\ \addlinespace[.8ex]

\textbf{Intervalle und Epsilon-Umgebungen}
\UebFormel{\begin{aligned}
[A,B]_{\mathbb R}&\coloneqq\{x\in\mathbb R\mid A\leq x\leq B\},\\
(A,B)_{\mathbb R}&\coloneqq\{x\in\mathbb R\mid A<x<B\},\\
B_\varepsilon(x)&\coloneqq
\{y\in\mathbb R\mid|y-x|<\varepsilon\},
\end{aligned}}
wobei \(\varepsilon>0\). Halboffene Intervalle und Strahlen werden
durch die entsprechenden nichtstrikten oder einseitigen Bedingungen definiert.
\UebQuellen{\FormulaRefAuto{RealBoundedIntervals}[def];
\FormulaRefAuto{RealUnboundedIntervals}[def];
\FormulaRefAuto{RealEpsilonNeighborhood}[def]}
&
\textit{Umgebungen sind offene Intervalle}
\UebFormel{B_\varepsilon(x)=(x-\varepsilon,x+\varepsilon)_{\mathbb R}.}
Die Ordnungs- und die Abstandsbeschreibung stimmen damit überein.
\UebQuellen{\FormulaRefAuto{RealEpsilonNeighborhoodInterval}[thm]}
\\ \addlinespace[.8ex]

\textbf{Archimedische Einbettung}
Für \(\eta\colon\mathbb N\to K\) in einer total geordneten Menge:
\UebFormel{\begin{aligned}
&\operatorname{Arch}(K,\preccurlyeq,\eta)\coloneqq\\
&\quad\forall x\in K\;\exists n\in\mathbb N\\
&\qquad(x\preccurlyeq\eta(n)\land x\neq\eta(n)).
\end{aligned}}
Im reellen Modell ist
\(\eta=\RealEmb\circ\RatEmbed\circ\IntEmbed\).
\UebQuellen{\FormulaRefAuto{ArchimedeanOrderedEmbedding}[def]}
&
\textit{Archimedizität, Ganzzahlabschnitt und Dichte}
\UebFormel{\begin{aligned}
&\forall x\in\mathbb R\;\exists n\in\mathbb N\;(x<\eta(n)),\\
&\forall x\in\mathbb R\;\exists z\in\mathbb Z:\\
&\quad\RealEmb(\RatEmbed(z))\leq x,\\
&\quad x<\RealEmb(\RatEmbed(z+\IntOne)),\\
&x<y\vdash\exists q\in\mathbb Q\;(x<\RealEmb(q)<y).
\end{aligned}}
\UebQuellen{\FormulaRefAuto{RealArchimedeanProperty}[thm];
\FormulaRefAuto{RealIntegerPartProperty}[thm];
\FormulaRefAuto{RationalDenseInReal}[thm]}
\\ \addlinespace[.8ex]

\textbf{Vollständig angeordneter Körper}
Die Strukturdefinition verbindet die Körpergesetze, eine kompatible
totale Ordnung mit \(0_K<1_K\) und die Supremumseigenschaft:
\UebFormel{\begin{aligned}
&\varnothing\neq S\subseteq K,\quad
\exists u\in K\;\forall s\in S\;(s\preccurlyeq u)\\
&\qquad\vdash\sup_K S\text{ existiert in }K.
\end{aligned}}
\UebQuellen{\FormulaRefAuto{CompleteOrderedFieldAxioms}[def]}
&
\textit{Modell und Eindeutigkeit}
Das Schnittmodell ist ein vollständig angeordneter Körper.
Für jeden solchen Körper \(K\) gibt es genau eine Bijektion
\(\Phi\colon\mathbb R\bij K\), die die folgenden Gesetze erfüllt:
\UebFormel{\begin{aligned}
&\Phi(x+y)=\Phi(x)\boxplus\Phi(y),\\
&\Phi(xy)=\Phi(x)\boxtimes\Phi(y),\\
&x\leq y\leftrightarrow\Phi(x)\preccurlyeq\Phi(y),\\
&\Phi(0)=0_K,\qquad\Phi(1)=1_K.
\end{aligned}}
\UebQuellen{\FormulaRefAuto{RealCutsCompleteOrderedField}[thm];
\FormulaRefAuto{CompleteOrderedFieldUniqueness}[thm]}
\\
\end{BandUebersicht}
